\chapter{Бимодульный допуск цепи когерентности}
\label{ch:v7-edge-coherence-bimodule-admission-gate}

\section{Что именно требуется вложить}

Спектральный родитель предыдущей главы использует цепь комплексных
размерностей $1\to6\to3$. Совпадение размерностей с отдельными
фермионными мультиплетами ещё не означает существование конечной
спектральной тройки. Для
\begin{equation}
 \mathcal A_{\rm SM}
 =\mathbb C\oplus\mathbb H\oplus M_3(\mathbb C)
 \label{eq:v7-coherence-bimodule-algebra}
\end{equation}
каждый неприводимый бимодуль имеет две координаты $(i,j)$, а его
комплексная размерность равна произведению размерностей простых модулей.

Для блока $D_{ij,kl}$ строгий первый порядок требует
\begin{equation}
 [[D,a],JbJ^{-1}]_{ij,kl}
 =(a_i-a_k)D_{ij,kl}(b_j-b_l)=0
 \quad\text{для всех }a,b.
 \label{eq:v7-coherence-bimodule-first-order}
\end{equation}
Следовательно, ненулевое ребро должно сохранять левую или правую
бимодульную координату.

\section{Неприводимый размерностный запрет}

Размерности простых комплексных модулей равны
\begin{equation}
 \dim\mathbb C=1,\qquad
 \dim_{\mathbb C}\mathbb H=2,\qquad
 \dim M_3=3.
 \label{eq:v7-coherence-bimodule-simple-dimensions}
\end{equation}
Поэтому одномерный неприводимый бимодуль имеет единственный тип
\begin{equation}
 \mathcal H_1^{(1)}=(\mathbb C,\mathbb C),
 \label{eq:v7-coherence-bimodule-one-dimensional}
\end{equation}
а шестимерные типы исчерпываются парой
\begin{equation}
 \mathcal H^{(6)}
 \in\{(\mathbb H,M_3),(M_3,\mathbb H)\}.
 \label{eq:v7-coherence-bimodule-six-dimensional}
\end{equation}
Одномерный тип не разделяет ни одной координаты ни с одним шестимерным
типом. Поэтому уже первое ребро любой неприводимой цепи $1\to6$ нарушает
\eqref{eq:v7-coherence-bimodule-first-order}.

\begin{theorem}[Неприводимый запрет цепи $1\to6\to3$]
Над неизменённой алгеброй $\mathcal A_{\rm SM}$ не существует
неприводимой бимодульной цепи размерностей $1\to6\to3$, удовлетворяющей
строгому первому порядку. Препятствие возникает на первом ребре и не
зависит от выбора трёхмерного конечного узла.
\end{theorem}

\section{Проверка фактически допущенного носителя}

В четырёхвершинном расширении единственный шестимерный физический узел есть
\begin{equation}
 Q_L=(\mathbb H,M_3,L),\qquad \dim_{\mathbb C}Q_L=6.
 \label{eq:v7-coherence-bimodule-q-node}
\end{equation}
Одномерные узлы нужной противоположной хиральности равны
\begin{equation}
 e_R,X_R=(\mathbb C,\mathbb C,R),
 \label{eq:v7-coherence-bimodule-one-nodes}
\end{equation}
а трёхмерные узлы равны
\begin{equation}
 u_R,d_R=(\mathbb C,M_3,R).
 \label{eq:v7-coherence-bimodule-three-nodes}
\end{equation}
Таким образом, имеются четыре размерностных кандидата
\begin{equation}
 \{e_R,X_R\}\longrightarrow Q_L
 \longrightarrow\{u_R,d_R\}.
 \label{eq:v7-coherence-bimodule-four-candidates}
\end{equation}
Второе ребро сохраняет правую координату $M_3$ и проходит первый порядок.
Первое меняет одновременно $\mathbb C\to\mathbb H$ и
$\mathbb C\to M_3$, поэтому запрещено. Real-сопряжение только переставляет
обе координаты и сохраняет тот же дефект.

\section{Почему прежнее поле всё же существовало}

Матрица $B\in M_{2\times3}(\mathbb C)$ была построена на пространстве
коэффициентов шести разрешённых стрелок полного графа:
\begin{equation}
 \{L_L,Y_L\}\times\{e_R,X_R,Y_R\}.
 \label{eq:v7-coherence-bimodule-edge-space}
\end{equation}
Это шестимерное пространство амплитуд рёбер, а не новый шестимерный
фермионный бимодуль. Аналогично,
\begin{equation}
 \Lambda^2\mathbb C^2\otimes\Lambda^2\mathbb C^3
 \simeq\mathbb C^3
 \label{eq:v7-coherence-bimodule-exterior-endpoint}
\end{equation}
является пространством составных миноров, а не уже существующей физической
вершиной. Поэтому спектральные следовые тождества предыдущего гейта
остаются верными на комплексе полей, но не превращаются автоматически в
фермионный конечный оператор.

\section{Минимальная приводимая альтернатива}

Первый порядок можно формально сохранить, если все три узла лежат в одном
бимодульном столбце:
\begin{equation}
 (\mathbb C,\mathbb C,R)
 \longrightarrow
 3(\mathbb H,\mathbb C,L)
 \longrightarrow
 3(\mathbb C,\mathbb C,R).
 \label{eq:v7-coherence-bimodule-reducible-chain}
\end{equation}
Размерности этой цепи действительно равны $1,6,3$. Однако крайние узлы
являются четырьмя независимыми копиями типа $(\mathbb C,\mathbb C,R)$,
а средний --- тремя копиями $(\mathbb H,\mathbb C,L)$. В допущенном
носителе имеются только
\begin{equation}
 n_{HCL}=2,\qquad n_{CCR}=2,
 \label{eq:v7-coherence-bimodule-available-multiplicities}
\end{equation}
тогда как цепь требует
\begin{equation}
 n_{HCL}=3,\qquad n_{CCR}=4.
 \label{eq:v7-coherence-bimodule-required-multiplicities}
\end{equation}
До аномального завершения пришлось бы добавить как минимум один левый
слабый дублет и два правых синглета. Простейшее вектороподобное завершение
добавляет также их три зеркальных партнёра. Кроме того, первый порядок
разрешает два независимых ребра, но сам не выводит нелинейную связь
$C_B=\frac12d(\Lambda^2)_B$ между ними.

\begin{nogocriterion}[Запрет переименования пространства стрелок]
Нельзя считать пространство коэффициентов Дираковских блоков физическим
фермионным модулем только из-за совпадения размерности шесть. Нельзя также
отождествлять внешний составной узел с трёхмерным кварковым модулем:
их бимодульные координаты и калибровочные действия различны.
\end{nogocriterion}

\section{Вердикт}

\begin{protocol}
\begin{enumerate}
 \item неприводимых одномерных типов: \textbf{только $(\mathbb C,\mathbb C)$};
 \item неприводимых шестимерных типов:
 \textbf{$(\mathbb H,M_3)$ и $(M_3,\mathbb H)$};
 \item допустимых неприводимых цепей $1\to6\to3$: \textbf{ноль};
 \item размерностных цепей в физическом носителе: \textbf{четыре};
 \item их допустимое первое ребро: \textbf{ни у одной};
 \item Real-сопряжение исправляет препятствие: \textbf{нет};
 \item пространство шести стрелок остаётся допустимым комплексом полей:
 \textbf{да};
 \item приводимая same-column цепь: \textbf{формально проходит};
 \item дефицит её физических вершин:
 \textbf{один дублет и два синглета};
 \item relation двух рёбер выводится первым порядком: \textbf{нет};
 \item итог:
 \textbf{неизменённый физический бимодульный носитель закрыт};
 \item следующий гейт:
 \textbf{вспомогательный BV/BRST или mapping-cone статус комплекса полей}.
\end{enumerate}
\end{protocol}

Машинный сертификат сохранён в
\path{s2t_v7_edge_coherence_bimodule_admission_gate_results.json}.

\begin{thebibliography}{9}
\bibitem{v7-coherence-bimodule-krajewski}
T.~Krajewski,
\emph{Classification of finite spectral triples},
Journal of Geometry and Physics 28 (1998), 1--30;
arXiv:hep-th/9701081.
\bibitem{v7-coherence-bimodule-vectorlike}
J.-H.~Jureit, C.~A.~Stephan,
\emph{On a Classification of Irreducible Almost-Commutative Geometries V},
Journal of Mathematical Physics 50 (2009), 072303; arXiv:0901.3214.
\end{thebibliography}